% (C) 2026 Ralf Stephan, in collaboration with Claude Code.
% Released under CC0 1.0 Universal (public-domain dedication).
% See https://creativecommons.org/publicdomain/zero/1.0/

\documentclass[11pt]{article}
\usepackage[T1]{fontenc}
\usepackage[utf8]{inputenc}
\usepackage{lmodern}
\usepackage[margin=2.5cm]{geometry}
\usepackage{amsmath,amssymb}
\usepackage{booktabs,tabularx}
% A results table longer than a page
\usepackage{xltabular}
\usepackage{braket}
\usepackage[hidelinks]{hyperref}

% Lean names and paths: verbatim, breakable at . _ /
\DeclareUrlCommand\lean{\urlstyle{tt}}
\def\UrlBreaks{\do\.\do\_\do\/}
\newcolumntype{L}{>{\raggedright\arraybackslash}X}
\newcolumntype{B}{>{\raggedright\arraybackslash}p{0.42\textwidth}}
\emergencystretch=6em

\title{Lean formalization of an Assumptions of Physics technical brief:\\ log}
\author{}
\date{}

\begin{document}
\maketitle

\noindent
This log records which statements of the technical brief in this directory were checked in Lean,
where the formal version differs from the brief, and any problems found in the brief.
The brief is available at \url{https://assumptionsofphysics.org/resources/briefs}.

\section*{Setup}

\begin{itemize}
\item Lean \texttt{v4.36.0-rc2}; Mathlib at commit \texttt{1b1ccf4b}; physlib (\texttt{master})
  and TauCeti (commit \texttt{e7e7f52}, on the same Mathlib) as further dependencies.
\item Build: \texttt{lake build} in the project root.
\item Every file is checked with \texttt{\#print axioms} on its main results. A result is
  recorded here only if it depends on no axioms beyond \lean{propext},
  \lean{Classical.choice} and \lean{Quot.sound}, and so has no \lean{sorry}.
\item Imports are kept to the Mathlib modules actually needed.
\item Material that belongs in Mathlib is kept in \lean{ForMathlib/}, in Mathlib's directory
  layout (library \lean{ForMathlib}), and imported by the brief files.
\end{itemize}


%======================================================================

\section{Unitary evolution as infinitesimal projective processes}

\paragraph{Files.}\leavevmode
\begin{itemize}
\item File: \lean{briefs/UnitaryProjectiveEvolution/UnitaryProjectiveEvolution.lean}
\item Library: \lean{UnitaryProjectiveEvolution}
\item Imports:\newline \lean{ForPhyslib.QuantumMechanics.VonNeumannResolvent},\newline
  \lean{ForPhyslib.QuantumMechanics.Wigner.AffineExtension},\newline
  \lean{ForMathlib.Analysis.SpecialFunctions.Exponential},\newline
  and the brief \lean{SchroedingerEvolution} (for \lean{IsEvolution}, \lean{DrUnit}, \lean{DrEv},
  \lean{DrEv.drUnit} and Proposition~1)
\item New file \lean{ForPhyslib/QuantumMechanics/VonNeumannResolvent.lean}: the von Neumann
  generator, the unitary channel of physlib's \lean{UnitaryOneParameterGroup}, the resolvent and
  its spectral identity, the second-order estimates, the largest outcome on a pure state with a
  Davis--Kahan bound, the resolvent exponential formula and time reversal. \lean{Pinching} gained
  the spectral decomposition \lean{sum_smul_spectralProj} and \lean{spectralProj_ne_zero}.
\item New \lean{ForMathlib} material:\newline
  \lean{ForMathlib/Analysis/InnerProductSpace/Trace.lean} (the operator norm of a positive
  operator is at most its trace).
\end{itemize}

\paragraph{Setting.} $E$ is a finite-dimensional complex Hilbert space, with density operators,
pure states and \lean{AffineOnDensities} as in the brief ``Multiple equivalent characterizations
of the Schr\"odinger equation''. $V$ is a norm-continuous unitary one-parameter group with
self-adjoint generator $K$, $V_t = e^{-itK}$; $\hbar = 1$, so $K = H/\hbar$.
\lean{vonNeumannGenerator}~$K$~$\rho = -i(K\rho - \rho K)$ is $L$,
\lean{unitaryChannel}~$V$~$t$~$\rho = V_t\rho V_t^\dagger$ is $T_t$, and
\lean{vonNeumannResolvent}~$V$~$\tau = (1 - \tau L)^{-1}$ is $\Psi_\tau$, a linear equivalence of
$\mathrm{End}(E)$. The definitions are
\begin{itemize}
\item \lean{IsProjectiveDecomposition}~$P$: star projections $P_a$, $P_aP_b = 0$ for
  $a \ne b$, $\sum_a P_a = 1$, over a finite index type;
\item \lean{IsStateAdaptiveProjective}~$\Psi$: $\Psi$ maps density operators to density
  operators, is affine on them, and for each density operator $\rho$ there is a projective
  decomposition with $\Psi\rho = \sum_a P_a\rho P_a$ (Definition~1);
\item \lean{IsOutcome}~$\Psi$~$P$~$p$~$Q$: $Q$ is pure and $F_aPF_a = pQ$ for one term of a
  projective decomposition with $\Psi P = \sum_a F_aPF_a$;
  \lean{IsBranch}~$\Psi$~$n$~$Q$~$p$: $Q_0$ is pure and each $Q_{k+1}$ is an outcome on $Q_k$
  of probability $p_k$, for $k < n$;
\item \lean{PerfectPrediction}~$\Psi$~$t$: from every pure $P$ there are, for all large $n$,
  branches $Q^{(n)}$ of $n$ steps of $\Psi_{t/n}$ with $Q^{(n)}_0 = P$, whose probability
  $\prod_k p^{(n)}_k$ tends to $1$ and whose final state $Q^{(n)}_n$ converges;
\item \lean{IsProjectiveLimit}~$\Phi$~$\Psi$: every $\Psi_\tau$ with $\tau > 0$ is
  state-adaptive projective, $(\Psi_\tau\rho - \rho)/\tau$ converges as $\tau \to 0^+$ for every
  density operator, $(\Psi_{t/n})^n\rho \to \Phi_t\rho$ for $t > 0$, and prediction is perfect
  for every $t > 0$;
\item \lean{DrProj}~$\Phi$ (DR-PROJ): $\Phi$ is a projective limit, and a second projective limit
  $\Phi'$ (reconstruction) satisfies $\Phi'_t(\Phi_tP) = P$ and $\Phi_t(\Phi'_tP) = P$ for pure
  $P$ and $t > 0$;
\item \lean{DrUnit} (DR-UNIT) and \lean{IsEvolution} are those of the brief on the
  Schr\"odinger equation.
\end{itemize}

\paragraph{Results.}
\begin{xltabular}{\textwidth}{@{}BL@{}}
\toprule
Brief & Lean \\
\midrule
\endfirsthead
\toprule
Brief & Lean \\
\midrule
\endhead
Definition~1 & \lean{IsStateAdaptiveProjective} \\
Proposition~2, $(I - \tau L)^{-1}$ exists & \lean{injective_vonNeumannStep},
  \lean{vonNeumannResolvent} \\
Eq.~(10) & \lean{vonNeumannResolvent_sub} \\
Eq.~(11) & \lean{spectralProj_mul_vonNeumannGenerator_mul} \\
Eq.~(12) & \lean{spectralProj_mul_mul_vonNeumannResolvent} \\
Eq.~(4), $\Psi_\tau\rho = \sum_a P_a^\rho\rho P_a^\rho$ & \lean{pinching_vonNeumannResolvent} \\
$\Psi_\tau$ is affine, positive, trace preserving & \lean{isDensity_vonNeumannResolvent},
  \lean{trace_vonNeumannResolvent} \\
Eq.~(7), and $\Psi_\tau$ is completely positive & \lean{vonNeumannResolvent_eq_laplaceResolvent},
  \lean{laplaceResolventCP} (any-dimension part) \\
Proposition~2 & \lean{isStateAdaptiveProjective_vonNeumannResolvent} \\
$\|L_K\| \le 2M/\hbar$ & \lean{norm_vonNeumannGenerator_apply_le} \\
Eq.~(16) & \lean{norm_vonNeumannResolvent_sub_sub_le}, \lean{norm_unitaryChannel_sub_sub_le} \\
Proposition~13, eq.~(14) & \lean{norm_vonNeumannResolvent_sub_unitaryChannel_le}:
  $\le 8\|K\|^2\tau^2$ \\
Proposition~13, eq.~(15), and $F_\tau(P)$ has rank one & \lean{exists_spectralProj_max_eq_proj}:
  $1 - q_\tau(P) \le 8\|K\|^2\tau^2$, $\|F_\tau(P) - T_\tau P\| \le 48\|K\|^2\tau^2$ \\
Proposition~17 & \lean{implemented_of_bijOnPure} \\
$(\Psi_\tau - I)/\tau = L\Psi_\tau \to L$ & \lean{tendsto_inv_smul_vonNeumannResolvent_sub} \\
Eq.~(20) & \lean{tendsto_vonNeumannResolvent_iterate}, error $\le 8\|K\|^2t^2/n$ \\
Eq.~(21), errors accumulate to $O(t^2/n)$ & \lean{perfectPrediction_vonNeumannResolvent},
  \lean{norm_sub_unitaryChannel_le_of_step} \\
``Replacing $H$ with $-H$'' & \lean{timeReversal}, \lean{unitaryChannel_timeReversal} \\
DR-UNIT $\Rightarrow$ DR-PROJ & \lean{DrUnit.drProj} \\
The branch has weight tending to one in the unconditioned output &
  \lean{isPositive_iterate_sub_of_isBranch}, \lean{norm_iterate_sub_le_of_isBranch} \\
$\Phi_t$ maps pure states to pure states & \lean{IsProjectiveLimit.isPure},
  \lean{isClosed_setOf_isPure} \\
$\Phi_t$ maps pure states bijectively onto pure states & \lean{DrProj.drEv} \\
Eq.~(27), phases & from the brief on the Schr\"odinger equation, \lean{DrEv.drUnit} \\
DR-PROJ $\Rightarrow$ DR-UNIT & \lean{DrProj.drUnit} \\
Theorem~18 & \lean{drUnit_iff_drProj} \\
A map affine on density operators is linear on them & \lean{exists_basis_isDensity},
  \lean{AffineOnDensities.basis_constr_apply}, \lean{AffineOnDensities.exists_linearMap} \\
$(T_{t/n})^n \to e^{tB}$ when $(T_\tau - I)/\tau \to B$ &
  \lean{tendsto_pow_exp_of_tendsto_smul_sub_one} \\
Affinity of the limit, eqs.~(28)--(29), composition for all durations &
  \lean{IsProjectiveLimit.exists_eq_exp} ($\Phi_t = e^{tB}$ on density operators),
  \lean{IsProjectiveLimit.isEvolution}, \lean{DrProj.isEvolution} \\
Theorem~18 without \lean{IsEvolution} & \lean{drUnit_iff_drProj_and_map_zero} \\
\bottomrule
\end{xltabular}

\paragraph{Differences from the brief.}
\begin{itemize}
\item Finite dimension only; the brief has a separable Hilbert space. Every generator is bounded,
  so the cutoffs $H_\tau = f_\tau(H)$, eqs.~(22)--(25) and the remark on unbounded generators are
  not needed and not formalized.
\item Norms are operator norms, not trace norms. Density operators have operator norm at most one,
  so the bounds of Proposition~13 are uniform over density operators as in the brief, with the
  explicit constants $8$ and $48$. In finite dimension trace-norm convergence follows, but a
  trace-norm version of (14) and (15) would carry the dimension.
\item Positivity and trace preservation of $\Psi_\tau$ come from eq.~(12), not from the integral
  representation~(7): $\Psi_\tau\rho = \sum_x P_x\rho P_x$, which is a density operator when
  $\rho$ is. The same identity gives injectivity of $1 - \tau L$, so Proposition~2 holds for
  every real $\tau$. Eqs.~(5)--(6) are not formalized. Eqs.~(7)--(9) and complete positivity come
  from the any-dimension part: there $\Psi_\tau$ is the integral~(7), and in finite dimension it
  equals $(1 - \tau L)^{-1}$ (\lean{vonNeumannResolvent_eq_laplaceResolvent}).
\item Eq.~(15) is proved for the largest eigenvalue, once $8\|K\|^2\tau^2 \le \frac12$; it is the
  only case the proof of Theorem~18 uses.
\item DR-PROJ is made precise as follows. Outcomes and branches refer to a projective
  decomposition that realizes $\Psi$ on the current pure state. Perfect prediction asks that the
  final branch state converges, without naming the limit; DR-PROJ $\Rightarrow$ DR-UNIT shows that
  it is $\Phi_tP$. Reconstruction is a second projective limit that inverts $\Phi_t$ on pure
  states. For the resolvents it is the limit of $(1 + \tau L)^{-1}$, as in the remark ``The
  finite process is not unitary''.
\item \lean{drUnit_iff_drProj} assumes \lean{IsEvolution} of the brief on the Schr\"odinger
  equation: $\Phi_t$ is affine and preserves density operators, $\Phi_0 = I$,
  $\Phi_{s+t} = \Phi_s\Phi_t$, and $t \mapsto \Phi_t\rho$ is continuous at $0^+$. The brief
  derives these from DR-PROJ (affinity of the limit, eqs.~(28)--(29), common refinements).
  \lean{IsProjectiveLimit.isEvolution} derives them too, but not by the brief's route. Each
  $\Psi_\tau$ agrees on density operators with a linear map $T_\tau$, because the density
  operators contain a basis of the operators; $(T_\tau - I)/\tau \to B$ in operator norm, because
  convergence on a basis is norm convergence; and $(T_{t/n})^n \to e^{tB}$ by telescoping. So
  $\Phi_t = e^{tB}$ on density operators for every $t > 0$, which gives affinity, preservation of
  density operators, continuity at $0^+$ and composition for all durations at once; eq.~(28) and
  the common refinements are not needed. DR-PROJ says nothing about $\Phi_0$, which is assumed to
  be the identity: \lean{drUnit_iff_drProj_and_map_zero} is DR-UNIT $\Leftrightarrow$ DR-PROJ
  with $\Phi_0 = I$.
\end{itemize}

\paragraph{Problems found in the brief.} No errors in the statements. Gaps and implicit points:
\begin{itemize}
\item Proposition~2 shows that the integral~(7) is a right inverse of $I - \tau L$. That it is
  $(I - \tau L)^{-1}$ also needs injectivity, which is not mentioned. In finite dimension it
  follows from eq.~(11): for self-adjoint $\sigma$, $\sigma = \sum_x P_x(\sigma - \tau L\sigma)P_x$
  over the spectral projectors of $\sigma$, so $(I - \tau L)\sigma = 0$ forces $\sigma = 0$.
\item Proposition~13 appeals to ``standard perturbation of an isolated eigenvalue'' without a
  reference or a gap condition. Here it is a Davis--Kahan bound. Let $q \ge \frac12$ be the
  largest eigenvalue of $\sigma = \Psi_\tau P$, $F = \ket{\phi}\bra{\phi}$ its spectral projector
  and $u = V_\tau\psi$. Then $\sigma - qF$ is positive with trace $1 - q$, so
  $\braket{w|\sigma w} \le (1-q)\|w\|^2$ for $w = u - \braket{\phi|u}\phi$. Comparing with
  $\sigma u = u + (\sigma - P_u)u$ gives $\|w\| \le 2\|\sigma - P_u\|$, and
  $\|F - P_u\| \le 3\|w\|$.
\item Section~3 states DR-PROJ only as ``the probabilities of the predicted and reconstructed
  state sequences tend to one''. The remark on prediction and reconstruction adds that the
  predicted states converge. This is not an extra assumption: since the branch is within
  $2(1-w)$ of $\Psi^nP \to \Phi_tP$, the predicted states converge to $\Phi_tP$ anyway. The
  finite-dimensional proof still uses the clause; the proof in any dimension does not
  (Section~\ref{sec:unitary-projective-any-dim}).
\item ``Perfect reconstruction gives the converse correspondence'' does not say how the
  reconstructed sequence relates to the predicted one. Injectivity and surjectivity of $\Phi_t$ on
  pure states need the reconstruction to invert $\Phi_t$; this is what DR-PROJ assumes here.
\item Composition for durations with irrational ratio, ``take common refinements with a
  remaining interval whose duration tends to zero'', is a gap. The remaining interval is
  harmless: with a common step $\tau$ the powers of $\Psi_\tau$ commute, so
  $\|\Psi_\tau^k(\Psi_\tau^j\rho - \rho)\|_1 \le j\|\Psi_\tau\rho - \rho\|_1
  \approx j\tau\|L\rho\|_1
  \to 0$ by the telescoping of eq.~(28). The problem is the step size. DR-PROJ defines $\Phi_t$
  only as $\lim_n (\Psi_{t/n})^n$. With $t/s$ irrational and $\tau = (t+s)/N$, the block of
  $k = \lfloor t/\tau\rfloor$ steps is $\Psi_\tau^k$, which is not of the form $(\Psi_{t/k})^k$,
  and DR-PROJ says nothing about its limit. Comparing $(\Psi_a)^k$ with $(\Psi_b)^k$ for nearby
  $a, b$ needs the $o(\tau)$ in eq.~(26) to be uniform along the states $\Psi_b^j\rho$, while
  eq.~(26) holds for each $\rho$ separately. Continuity at zero does not supply this. In finite
  dimension $L$ is bounded, eq.~(26) holds in operator norm, and telescoping gives
  $\|(\Psi_{t/n})^n - e^{tL}\| \le n\,o(t/n)\,e^{t\|L\|} \to 0$, so $\Phi_t = e^{tL}$ and
  composition follow at once (formalized, \lean{IsProjectiveLimit.exists_eq_exp}). In infinite
  dimension Chernoff's product formula gives the same,
  uniformly on compact $t$-intervals, provided the closure of $L$ generates a contraction
  semigroup; the brief's ``single, possibly unbounded generator $L$'' must be read to include
  this, or it must be assumed.
\item ``The phases of the implementing unitaries may be chosen'' is the lifting step already noted
  for the brief on the Schr\"odinger equation; it is filled in the same way.
\item The remarks ``The qubit picture'' ($\Psi_\tau = \frac12(I + C_\tau)$ with the Cayley
  transform $C_\tau$) and ``The finite process is not unitary'' are not formalized.
\end{itemize}

\subsection{Hilbert spaces of any dimension}\label{sec:unitary-projective-any-dim}

This subsection covers Definition~1, Propositions~2 and~13, eq.~(20) and Theorem~18 in both
directions, with an unbounded generator.

\paragraph{Files.}\leavevmode
\begin{itemize}
\item File: \lean{briefs/UnitaryProjectiveEvolution/UnitaryProjectiveEvolution/InfiniteDimensional.lean}
\item Library: \lean{UnitaryProjectiveEvolution}, which now builds every module
  \lean{UnitaryProjectiveEvolution.*}
\item Imports:\newline \lean{ForPhyslib.QuantumMechanics.ForwardHamiltonian},\newline
  \lean{ForPhyslib.QuantumMechanics.SpectralCutoff},\newline
  \lean{ForPhyslib.QuantumMechanics.VonNeumannResolvent.Laplace},\newline
  \lean{Mathlib.NumberTheory.Real.Irrational} (for \lean{Irrational}),\newline
  and \lean{SchroedingerEvolution.InfiniteDimensional} (for \lean{IsEvolution}, \lean{DrUnit},
  \lean{DrEv} and \lean{DrEv.drUnit}; it brings \lean{StronglyContinuousUnitaryGroup} and
  \lean{AffineOnDensityOperators})
\item New file \lean{ForPhyslib/QuantumMechanics/VonNeumannResolvent/Laplace.lean}: the resolvent
  of a bounded generator in any dimension as the Laplace integral of the unitary channel, its
  inverse identities, complete positivity, preservation of density operators, its agreement with
  $(1 - \tau L)^{-1}$ in finite dimension, the eigenbasis decomposition, the
  second-order estimates and the largest outcome in trace norm, the exponential formula and the
  generator limit.
\item New file \lean{ForPhyslib/QuantumMechanics/SpectralCutoff.lean}: the cutoff groups
  $\int e^{itc_M(x)}\,d\mu_S(x)$ through PhyslibAlpha's bounded spectral integral, the bound
  $\|H_M\| \le M$ on their generators, strong convergence as $M \to \infty$, and the spectral
  projectors $P_m = \mu_S[-m,m]$.
\item Additions to \lean{ForPhyslib/QuantumMechanics/TraceNorm.lean}: \lean{traceENorm}, the trace
  norm $\operatorname{tr}|T| \in [0,\infty]$ of any operator, with its calculus; the trace-norm
  distance of two pure states (\lean{traceENorm_proj_sub_proj_le}); extension of trace-norm bounds
  from pure states to density operators (\lean{traceENorm_apply_le_of_forall_proj}); trace-norm
  convergence of $U_a\rho U_a^\dagger$ when $U_a \to U$ strongly (\lean{tendsto_traceENorm_conj});
  density of the density operators supported in projections $P_a \to 1$
  (\lean{exists_isDensityOperator_mul_eq}); closedness of the pure states and of the density
  operators under trace-norm limits (\lean{isPure_of_tendsto_traceENorm} and
  \lean{isDensityOperator_of_tendsto_traceENorm}); uniqueness of trace-norm limits; and the
  affinity, contraction and telescoping bounds for iterates of affine maps.
\item Addition to \lean{ForPhyslib/QuantumMechanics/Stone.lean}:\newline
  \lean{StronglyContinuousUnitaryGroup.exists_expIntegral}, every strongly continuous unitary
  group is $t \mapsto \int e^{itx}\,d\mu_S(x)$ for a projection-valued measure $\mu_S$ on
  $\mathbb{R}$.
\item New \lean{ForMathlib} material:\newline
  \lean{ForMathlib/Analysis/CStarAlgebra/CompletelyPositiveMap.lean} (an integral of completely
  positive maps is completely positive, \lean{CompletelyPositiveMap.integral}).
\end{itemize}

\paragraph{Setting.} $E$ is a complex Hilbert space of any dimension; separability is not used.
Density operators (\lean{IsDensityOperator}) and pure states are as in
the section ``Hilbert spaces of any dimension'' of
\lean{briefs/SchroedingerEvolution/formalization-log.tex}, and DR-UNIT is \lean{DrUnit} from there:
$\Phi_t\rho = U_t\rho U_t^\dagger$ for a strongly continuous unitary group, whose generator may be
unbounded. Limits are in the trace norm, through \lean{traceENorm}. For a norm-continuous group
$V_t = e^{-itK}$ with bounded generator, $L$, $T_t$ and $\hbar = 1$ are as in finite dimension, and
\lean{laplaceResolvent}~$V$~$\tau = \int_0^\infty e^{-s}T_{\tau s}\,ds$ is $\Psi_\tau$, a bounded
operator on $\mathrm{End}(E)$. The definitions are those of finite dimension, with these changes:
\begin{itemize}
\item \lean{IsProjectiveDecomposition}~$P$: the index type is arbitrary and $\sum_a P_a = 1$
  converges strongly; \lean{IsStateAdaptiveProjective}, \lean{IsOutcome}: the sums
  $\sum_a P_a\rho P_a$ converge strongly;
\item \lean{PerfectPrediction}: the final branch state converges in trace norm;
\item \lean{IsProjectiveLimit}~$\Phi$~$\Psi$: $(\Psi_{t/n})^n\rho \to \Phi_t\rho$ in trace norm,
  and the infinitesimal condition is eq.~(26): there is a set $D$ of density operators, dense in
  trace norm, such that $(\Psi_\tau\rho - \rho)/\tau$ converges in trace norm as $\tau \to 0^+$ for
  every $\rho \in D$.
\end{itemize}

\paragraph{Results.}
\begin{xltabular}{\textwidth}{@{}BL@{}}
\toprule
Brief & Lean \\
\midrule
\endfirsthead
\toprule
Brief & Lean \\
\midrule
\endhead
Definition~1 & \lean{IsStateAdaptiveProjective}, \lean{IsProjectiveDecomposition} \\
Eq.~(7), $\Psi_\tau$ as a mixture of unitary channels & \lean{laplaceResolvent} \\
Eqs.~(8)--(9), $(I - \tau L)\Psi_\tau = I$ & \lean{integral_laplace_vonNeumannStep},
  \lean{laplaceResolvent_sub} \\
$\Psi_\tau(I - \tau L) = I$ & \lean{laplaceResolvent_vonNeumannStep} \\
$\Psi_\tau$ is affine, positive, trace preserving & \lean{laplaceResolvent_nonneg},
  \lean{isDensityOperator_laplaceResolvent} \\
Proposition~2, $\Psi_\tau$ is completely positive & \lean{laplaceResolventCP},
  \lean{coe_laplaceResolventCP} \\
Eq.~(11), on an eigenbasis & \lean{inner_vonNeumannGenerator_apply_eq_zero} \\
Eqs.~(4), (12) & \lean{exists_hilbertBasis_laplaceResolvent}: $\Psi_\tau\rho = \sum_i
  P_i\rho P_i$ with $P_i\rho P_i = \mu_iP_i$ \\
Proposition~2 & \lean{isStateAdaptiveProjective_laplaceResolvent} \\
$\|L_K\| \le 2M/\hbar$ on trace-class operators & \lean{traceENorm_vonNeumannGenerator_le} \\
Proposition~13, eq.~(14), in trace norm &
  \lean{traceENorm_laplaceResolvent_sub_unitaryChannel_le}: $\le 80\|K\|^2\tau^2$ \\
Proposition~13, eq.~(15), in trace norm & \lean{exists_max_laplaceResolvent_proj}:
  $1 - q_\tau(P) \le 8\|K\|^2\tau^2$, $\|F_\tau(P) - T_\tau P\|_1 \le 64\|K\|^2\tau^2$ \\
$(\Psi_\tau - I)/\tau = L\Psi_\tau \to L$, bounded $K$ &
  \lean{tendsto_inv_smul_laplaceResolvent_sub}, error $\le 4\|K\|^2\tau$ \\
Eq.~(20), bounded $K$ & \lean{tendsto_laplaceResolvent_iterate},
  \lean{traceENorm_laplaceResolvent_iterate_sub_le}: error $\le 80\|K\|^2t^2/n$ \\
Eq.~(22) & \lean{cutoffLevel}: $M_\tau = \tau^{-1/4}$, \lean{mul_cutoffLevel_sq},
  \lean{tendsto_cutoffLevel} \\
$H_\tau = f_\tau(H)$, $\|H_\tau\| \le M_\tau$ & \lean{cutoffGroup},
  \lean{norm_cutoffGroup_generator_le} \\
$e^{tH_\tau/(i\hbar)} \to e^{tH/(i\hbar)}$ strongly & \lean{tendsto_cutoffGroup_apply} \\
Eq.~(23) & \lean{cutoffResolvent} \\
Eq.~(26), on the core of bounded spectral support & \lean{infinitesimal_cutoffResolvent},
  \lean{laplaceResolvent_cutoffGroup_eq}, \lean{cutoffGroup_mul_spectralCutoffProj} \\
The generator of family~(23) is $[H,\rho]/(i\hbar)$, on the core &
  \lean{cutoffResolvent_generator}, \lean{expGroup_mul_spectralCutoffProj},
  \lean{exists_mem_domain_of_spectralCutoffProj}, \lean{forwardHamiltonian_expGroup} \\
The core is dense & \lean{exists_isDensityOperator_mul_eq},
  \lean{tendsto_spectralCutoffProj_apply} \\
Eqs.~(24)--(25), hence eq.~(20) for unbounded $H$ & \lean{tendsto_cutoffResolvent_iterate},
  \lean{traceENorm_sub_unitaryChannel_le_of_step}, \lean{tendsto_traceENorm_conj} \\
Perfect prediction, unbounded $H$, total failure $O(t\tau M_\tau^2)$ &
  \lean{perfectPrediction_cutoffResolvent} \\
$U_t = e^{tH/(i\hbar)}$ for a spectral measure & \lean{StronglyContinuousUnitaryGroup.exists_expIntegral} \\
``Replacing $H_\tau$ with $-H_\tau$'' & \lean{timeReversal} \\
DR-UNIT $\Rightarrow$ DR-PROJ & \lean{DrUnit.drProj}, from \lean{isProjectiveLimit_cutoffResolvent}
  and \lean{exists_isProjectiveLimit} \\
The branch has weight tending to one in the unconditioned output &
  \lean{IsOutcome.nonneg}, \lean{nonneg_iterate_sub_of_isBranch},
  \lean{traceENorm_iterate_sub_le_of_isBranch}: $\|\Psi^n Q_0 - Q_n\|_1 \le 2(1 - w)$ \\
$\Phi_t$ maps pure states to pure states & \lean{IsProjectiveLimit.isPure},
  \lean{isPure_of_tendsto_traceENorm} \\
$\Phi_t$ maps pure states bijectively onto pure states & \lean{DrProj.drEv} \\
Proposition~17 and eq.~(27), phases & from the brief on the Schr\"odinger equation,
  \lean{DrEv.drUnit} (the section ``Hilbert spaces of any dimension'' of
\lean{briefs/SchroedingerEvolution/formalization-log.tex}) \\
DR-PROJ $\Rightarrow$ DR-UNIT & \lean{DrProj.drUnit} \\
Theorem~18 & \lean{drUnit_iff_drProj} \\
Every finite repetition is affine; $\Phi_t$ is affine and keeps density operators &
  \lean{AffineOnDensityOperators.iterate}, \lean{IsProjectiveLimit.affine},
  \lean{IsProjectiveLimit.isDensityOperator} \\
$\Psi_\tau$ and $\Phi_t$ contract the trace norm &
  \lean{AffineOnDensityOperators.traceENorm_sub_le}, \lean{IsProjectiveLimit.traceENorm_sub_le} \\
Eq.~(28), telescoping & \lean{IsProjectiveLimit.traceENorm_sub_le_mul},
  \lean{AffineOnDensityOperators.traceENorm_iterate_sub_le} \\
Continuity at zero, extended by contractivity & \lean{IsProjectiveLimit.tendsto_traceENorm_sub},
  \lean{IsProjectiveLimit.continuousWithinAt} \\
Eq.~(29), composition for durations of rational ratio & \lean{IsProjectiveLimit.map_add_mul},
  \lean{IsProjectiveLimit.map_add_of_not_irrational} \\
The evolution structure from DR-PROJ & \lean{IsProjectiveLimit.isEvolution}, assuming
  $\Phi_0 = I$ and composition for durations of irrational ratio \\
Theorem~18 without \lean{IsEvolution} & \lean{drUnit_iff_drProj_and_map_zero_and_map_add} \\
\bottomrule
\end{xltabular}

\paragraph{Differences from the brief.}
\begin{itemize}
\item Limits and the bounds of Proposition~13 are in the trace norm, as the brief's remark on
  prediction and reconstruction asks. \lean{traceENorm} is $\operatorname{tr}|T|$ in $[0,\infty]$
  for every operator, so convergence statements carry no trace-class witnesses; it is finite
  exactly on the trace class, where it is PhyslibAlpha's trace norm.
\item In Definition~1 the projective decomposition has an arbitrary index type and its sums
  converge strongly. For a density operator $\rho$ the sum $\sum_a P_a\rho P_a$ converges in trace
  norm anyway, so nothing is lost; in a separable space at most countably many $P_a$ are nonzero.
\item $\Psi_\tau$ is defined by the integral~(7), which converges in operator norm on
  $\mathrm{End}(E)$ for every real $\tau$. It is a two-sided inverse of $I - \tau L$: the brief's
  eqs.~(8)--(9) give the right inverse, and since $L$ commutes with every $T_s$ the same integral
  is a left inverse. Positivity comes from the integral. That $\Psi_\tau\rho$ is trace class with
  trace one uses the integral again, now in the trace class $\mathcal{S}_1$, where
  $s \mapsto T_s\rho$ is continuous and the trace is a continuous functional. Interchanging a
  basis sum with the integral would need a countable basis.
\item The decomposition of Proposition~2 uses the rank-one projectors $P_i$ onto an eigenbasis of
  $\sigma = \Psi_\tau\rho$, not the spectral projectors that group equal eigenvalues. Both satisfy
  Definition~1. Eq.~(11) becomes $\braket{b_i|L\sigma\,b_i} = 0$, so
  $P_i\rho P_i = \braket{b_i|\rho\,b_i}P_i = \mu_iP_i$, and ``every $P_a^P$ with $\lambda_a > 0$ has
  rank one'' is automatic. For a pure state the largest eigenvalue is at least
  $1 - 8\|K\|^2\tau^2 > \frac12$ and the eigenvalues sum to one, so it is simple and its
  eigenprojector is the brief's $F_\tau(P)$.
\item Proposition~13 is proved through the finite-dimensional operator-norm argument: the
  Davis--Kahan bound $\|w\| \le 2\|\sigma - P_u\|$ for the component $w$ of $u = V_\tau\psi$
  orthogonal to the top eigenvector $\phi$, then $\|P_\phi - P_u\|_1 \le 4\|w\|$ for unit vectors.
  For a pure state $P$, $\Psi_\tau P - qF_\tau(P)$ is positive with trace $1-q$, which gives
  eq.~(14) on pure states with constant $80$. It extends to all density operators by summing the
  canonical expansion in $\mathcal{S}_1$. For $8\|K\|^2\tau^2 > \frac12$ the bound is trivial,
  since $\|\Psi_\tau\rho - T_\tau\rho\|_1 \le 2$, so eq.~(14) holds for every $\tau \ge 0$.
\item The cutoff of eq.~(22) is $M_\tau = \tau^{-1/4}$, so $\tau M_\tau^2 = \sqrt\tau$.
  $H_\tau$ is never formed as an operator of its own. The cutoff group
  $W_t = \int e^{itc_M(x)}\,d\mu_S(x)$, with $c_M(x) = \max(-M,\min(x,M))$, is a norm-continuous
  unitary group because $\|W_t - W_s\| \le M|t-s|$, and $H_\tau$ is its bounded generator, of norm
  at most $M$ by the same Lipschitz bound. Up to the sign convention of physlib's
  \lean{UnitaryOneParameterGroup}, it is the brief's $f_\tau(H)$. Strong convergence of the cutoff
  groups is dominated convergence in
  $\|(W_t - U_t)x\|^2 = \int |e^{itc_M(x)} - e^{itx}|^2\,d\mu_x$.
\item In \lean{IsProjectiveLimit}, eq.~(26) does not name the generator: $(\Psi_\tau\rho -
  \rho)/\tau$ converges for $\rho$ in a dense set $D$ of density operators. $D$ is the brief's
  core, the density operators with $P_m\rho = \rho$ for $P_m = \mu_S[-m,m]$ and some $m$. On
  such $\rho$,
  $W^{(M)}_tP_m = W^{(m)}_tP_m$ for $M \ge m$, so once $M_\tau \ge m$, $\Psi_\tau\rho$ is the
  resolvent of the fixed bounded group $W^{(m)}$, and the bounded case applies. The core is dense
  because $P_m \to 1$ strongly: $P_m\rho P_m/\operatorname{tr}(P_m\rho P_m) \to \rho$ in trace norm.
  The limit is $L_{K_m}\rho$, the commutator with the bounded part $K_m$ of $H$. It is the
  brief's $[H,\rho]/(i\hbar)$ in the following sense (\lean{cutoffResolvent_generator}): for the
  self-adjoint $H$ with $U_t = e^{-itH}$ (Stone's theorem), $\rho$ maps into the domain of $H$,
  and $L\rho\,\psi = -i(H\rho\psi - \rho H\psi)$ for $\psi$ in the domain, which determines the
  bounded $L\rho$. On the range of $P_m$, $U_t = W^{(m)}_t$, so $H = K_m$ there; this uses that
  $H$ is the forward Hamiltonian of $U$ from the brief on the Schr\"odinger equation
  (\lean{forwardHamiltonian_expGroup}). Then $\rho H\psi = \rho K_m\psi$ by symmetry of $H$.
\item DR-UNIT $\Rightarrow$ DR-PROJ needs no \lean{IsEvolution}, as in finite dimension.
  Reconstruction uses the cutoff resolvents of the reversed group $t \mapsto U_{-t}$, which has its
  own spectral measure.
\item DR-PROJ $\Rightarrow$ DR-UNIT is proved over \lean{IsEvolution} (\lean{DrProj.drUnit}) and
  over two of its clauses only (\lean{drUnit_iff_drProj_and_map_zero_and_map_add}). The
  finite-dimensional derivation of \lean{IsEvolution}, through $\Phi_t = e^{tB}$, needs a finite
  basis of density operators; \lean{IsProjectiveLimit.isEvolution} follows the brief instead.
  Each $\Phi_t$ is a trace-norm limit of affine maps that keep density operators, so it is one
  too, and contracts the trace norm (Simon's lemma, the brief's ``extend to positive
  trace-preserving linear maps''). Eq.~(28) is the telescoping bound
  $\|\Psi^n\rho - \rho\|_1 \le n\|\Psi\rho - \rho\|_1$ on the dense set of eq.~(26), and
  contractivity extends continuity at $0^+$ to every density operator. Eq.~(29) is proved as
  $\Phi_{(p+q)u} = \Phi_{pu}\circ\Phi_{qu}$, along $n = (p+q)m$, which gives composition whenever
  $s/t$ is rational. Two clauses remain assumptions: $\Phi_0 = I$, since \lean{IsProjectiveLimit}
  defines $\Phi_t$ only for $t > 0$ (the brief obtains it from $\Psi_0 = I$), and
  $\Phi_{s+t} = \Phi_s\circ\Phi_t$ for $s/t$ irrational, the gap noted in the finite-dimensional
  part.
\item The branch estimate is the finite-dimensional one, with traces in $[0,\infty]$: along a
  branch of probability $w$,
  $\Psi^nQ_0 - wQ_n$ is positive with trace $1 - w$, so it is $(1-w)R$ for a density operator $R$,
  and $\|\Psi^nQ_0 - Q_n\|_1 \le 2(1-w)$. In finite dimension the pure states are closed because
  they are the image of the compact unit sphere; here they are closed in $\mathcal{S}_1$, as the
  positive idempotents of trace norm one. Proposition~17 is Proposition~1 of the brief on the
  Schr\"odinger equation in any dimension, which needs affinity under all convex combinations;
  \lean{IsStateAdaptiveProjective} and \lean{IsEvolution} both assume it.
\item The proof of DR-PROJ $\Rightarrow$ DR-UNIT does not use the clause of
  \lean{PerfectPrediction} that the final branch states converge. The bound
  $\|\Psi^nP - Q_n\|_1 \le 2(1-w_n)$ and $\Psi^nP \to \Phi_tP$ already give $Q_n \to \Phi_tP$.
  The clause is kept so that the definitions match those of finite dimension.
\item Complete positivity is Mathlib's \lean{CompletelyPositiveMap}: $\Psi_\tau \otimes
  \mathrm{id}$ is positive on the $k \times k$ matrices of operators for every $k$. It holds
  because every unitary channel is a $*$-automorphism, and an integral of completely positive
  maps is completely positive, since the positive operators form a closed cone. The remark on the
  resolvent of an unbounded generator is not formalized.
\end{itemize}

\paragraph{Problems found in the brief.} No errors. Gaps beyond those of the finite-dimensional
part:
\begin{itemize}
\item Proposition~2 uses the spectral decomposition $\sigma = \sum_a\lambda_aP_a^\rho$ of
  $\sigma = \Psi_\tau(\rho)$. In infinite dimension this needs $\sigma$ to be compact, which
  follows from $\sigma$ being trace class; the brief does not show that it is. The decomposition
  may have infinitely many terms and includes the kernel of $\sigma$, and its sum converges
  strongly.
\item The core of eq.~(26) is asserted to be dense without proof. Density follows from $P_m \to 1$
  strongly and trace-norm continuity of conjugation.
\item Composition for durations of irrational ratio is the gap noted in the finite-dimensional
  part. Here no telescoping estimate closes it: $L$ may be unbounded, and eq.~(26) holds only
  for each $\rho$ of a dense set, not uniformly.
\item Eq.~(25), ``strong convergence implies convergence of the corresponding conjugations on
  every trace-class operator'', is standard, but not obvious: for density operators it follows from
  the canonical expansion and the rank-one estimate $\|P_x - P_y\|_1 \le 2\|x - y\|$.
\end{itemize}

\end{document}
